\documentclass[10pt]{article}
\usepackage[margin=20mm]{geometry}
\usepackage{amsmath,amssymb,mathtools,fontspec,unicode-math,xurl,hyperref}
\IfFontExistsTF{DejaVu Serif}{\setmainfont{DejaVu Serif}}{\setmainfont{Latin Modern Roman}}
\setmathfont{latinmodern-math.otf}
\setlength{\parindent}{0pt}\setlength{\parskip}{6pt}\sloppy
\title{Initial scientific review
A quadratically enriched correspondence theorem}
\author{AI-assisted technical review}
\date{}
\begin{document}
\maketitle
Assignment. Authors. Andrés Jaramillo Puentes and Sabrina Pauli. Exact source. arXiv:2309.11706v2, submitted 13 April 2024, 02:24:30 UTC; 44 pages. The journal identifier is 10.4171/JEMS/1694 (https://doi.org/10.4171/JEMS/1694), published online on 15 September 2025. This report concerns the frozen v2; equivalence with the later journal version is not assumed.\par
\section*{Assessment and scope}
Major revision: two defects in the reconstruction argument, with two additional background corrections. The stated characteristic bound does not ensure the separability and nodality used by the local construction. Independently, the deformation-pattern algebra loses a coefficient-dependent Kummer twist already in characteristic zero. Both affect the actual reconstruction and its fields of definition. The examples below do not establish a contradiction to every version of the final aggregate formula: the characteristic-zero example has the same total transferred weight after correcting its algebra.\par
All 44 pages, including references, were read; 32 selected source pages were inspected as images. The review traced the node norm, trace forms, triangle and parallelogram contributions, deformation patterns, marked-edge extensions, the dimension product, Witt-class cancellation and the passage to del Pezzo invariance. Characteristic, separability, simplicity, boundary edge lengths, genus and genericity restrictions were kept distinct.\par
\section*{Confirmed problems}
Problem 1 (Major: the diameter bound does not protect the local construction).  The characteristic assumption occurs in Theorem 1.6, p. 7, the notation on p. 9, Theorem 4.1, p. 21, and Theorem 6.2, pp. 32–33. Lemmas 5.5–5.6, pp. 24–26, and Remark 6.5, p. 34, use it to obtain distinct roots, nodal curves and finite étale algebras. The original polygon’s ordinary Euclidean diameter cannot supply these conclusions. No different definition of diameter was located in the source.\par
Work over \(k=\overline{\mathbb F}_{11}\), and put \[\Delta=\operatorname{Conv}\{(0,0),(3,1),(1,4)\}.\] Its squared diameter is \(17\), its double area is \(11\), and its three edges are primitive. Thus it satisfies \(11>\operatorname{diam}(\Delta)\); Pick’s formula gives five interior lattice points. The affine unimodular map \[T(a,b)=(a-2b+7,-a+3b)\] has determinant \(1\) and sends its vertices to \((7,0),(8,0),(0,11)\). In the notation of Lemmas 5.5–5.6, \[m=11,\qquad p=7,\qquad q=8,\qquad d_p=d_q=1.\] Take \(\alpha=\beta=1\). The boundary constants \(a=-1,\ b=\varepsilon_1=\varepsilon_2=1\) satisfy (14), and the parameterization (13) becomes \[x=t^{11},\qquad y=t^7(t-1),\qquad
 F(x,y)=y^{11}-x^8+x^7=0.                         \tag{A}\] All three boundary intersections are single, smooth branches: their contact orders are the three primitive edge lengths, namely one. The vertex coefficients are nonzero, so the curve misses the toric fixed points.\par
This curve is irreducible and rational. Indeed \(x^7(x-1)\) is not an eleventh power in \(k(x)\), while \([k(t):k(t^{11})]=11\) and \(d(t^7(t-1))/dt=t^6(8t-7)\ne0\). Consequently \(k(t^{11},t^7(t-1))=k(t)\), so the smooth projective normalization is \(\mathbb P^1\).\par
Nevertheless (A) is not nodal. Since \[F_y=0,\qquad F_x=x^6(7-8x),\] its unique torus singular point is \((5,1)\). In local coordinates \(X=x-5,\ Y=y-1\), \[F(5+X,1+Y)=Y^{11}+2X^2+6X^3+\cdots .\] Its tangent cone is the double line \(2X^2\), not two distinct lines. Equivalently, the normalization is pointwise injective because \(t_1^{11}=t_2^{11}\) implies \(t_1=t_2\), and it fails to be immersive at \(t=5\). In the original coordinates the same example is \(1-U^3V+UV^4=0\), under \(U=x/y,\ V=y^3/x^2\); thus the singularity is not an artefact of changing coordinates.\par
The algebra (15) is now \[F'=k[\beta]/(\beta^{11}-1)
    =k[\beta]/(\beta-1)^{11},\] which is nonreduced and not étale. Its multiplication trace is zero, so it cannot provide the nondegenerate trace form used in Proposition 2.8. The claimed weight in Lemma 5.6 also contains a zero power of \(m=11\), rather than a unit defining a rank-one form. These are actual failures of the core nodality and trace construction, not a numerical tolerance issue. The polynomial identities, local expansion and lattice data were independently computed exactly modulo \(11\).\par
The dependency is direct: Lemmas 5.5–5.6 feed (42), the algebra \(F\), then \(M\) in (46), and the trace calculation in Section 6.2. The counterexample does not satisfy the additional nodality assumption imposed on already chosen component curves in Theorem 4.1; it refutes the assertion that the subsequent recipe always supplies such components under the advertised characteristic condition. A full generic-marked-point counterexample to the final enumerative equality has not been constructed here.\par
The primary preprint of Tyomkin’s cited work, arXiv:1001.1554v4, Theorem 6.2 on pp. 46–47, was checked. It explicitly assumes regularity over the residue field and characteristic-prime edge multiplicities; it does not state the diameter criterion attributed to it here. Its theorem and proof therefore cannot simply replace those checks.\par
A repair must control the relevant lattice indices and local exponents, and verify the regularity needed by patchworking. Unimodular transformations do not preserve Euclidean diameter. Restricting to characteristic zero excludes this example; a complete corrected positive-characteristic range is not proved in this review.\par
Problem 2 (Major: the deformation-pattern algebra drops its prescribed coefficient).  Definition 3.21, p. 20, fixes the two sloping boundary truncations of a deformation pattern, hence its coefficient of \(x^my\). Proposition 5.10(2), p. 30, identifies the algebra of choices, for odd \(m\), with \(F[u]/(u^m-1)\). Equation (33) silently fixes a special leading coefficient. Equations (43)–(44), pp. 37–38, then use this algebra for the general reconstruction.\par
The prescribed coefficient can arise from compatible component curves, rather than arbitrary boundary data. Over \(\mathbb Q\), take the adjacent triangles \[\Delta_\pm=\operatorname{Conv}\{(0,0),(3,0),(0,\pm1)\}\] and component equations \(2(x-1)^3+y\) and \(2(x-1)^3+y^{-1}\), respectively. Both curves are smooth and rational. Their common-edge truncation is \(2(x-1)^3\), giving contact of order three at \(x=1\); the other edges are primitive. The shift \(x=X+1\) in Section 3.3, pp. 18–20, yields the fixed sloping truncations \(y^2+2X^3y\) and \(2X^3y+1\) after multiplying the refinement by \(y\). Thus the following data satisfy the compatibility requirement of Definition 3.21.\par
Take \(F=k=\mathbb Q\), \(a=b=c=1\), \(m=3\), and prescribed leading coefficient \(a_3=2\). Set \[G=y^2+y(2x^3+Bx+C)+1.\] The missing \(x^2y\) coefficient is zero, as required. Two distinct singularities require the two critical values of \(2x^3+Bx+C\) to be \(2\) and \(-2\). If the critical points are \(\pm r\), their equations give \[B=-6r^2,\qquad C=0,\qquad B^3=-54.\] Conversely these equations suffice. Write \(B=-3u\), \(u^3=2\). An exact factorization is \[(2x^3-3ux)^2-4
       =4(x^2-u/2)^2(x^2-2u).                      \tag{B}\] There are exactly two ordinary nodes, \[x^2=u/2,\qquad y=ux,\] and at either one \(-\det\operatorname{Hess}(G)=-12u^2\ne0\). Removing the squared factor in (B) gives the smooth conic \(w^2=x^2-2u\); this proves geometric rationality. The boundary truncations are fixed. The two sloping edges are primitive, and the vertical truncation \(y^2+1\) has distinct roots, so there are no further boundary singularities.\par
The algebra of coefficient choices is therefore \[\mathbb Q[B]/(B^3+54)
       \cong\mathbb Q[u]/(u^3-2),\] a field by Eisenstein’s criterion. It is not \(\mathbb Q[u]/(u^3-1)\), which has a \(\mathbb Q\)-factor. In particular the printed algebra predicts a rational coefficient choice which does not exist. Formula (33), with these \(a,b,c\), instead has leading coefficient \(8\), not the prescribed \(2\). Changing the torus coordinate to obtain that normalization requires a cube root; descent retains precisely the twist that was omitted.\par
This damages the construction of a curve over the specific algebra \(M\) in Section 6.1: the wrong coefficient algebra cannot index the compatible patterns with their actual fields of definition. That is a central arithmetic reconstruction claim. The Major rating concerns this proved failure, not a claim that the final weighted count has been disproved. Indeed, for this example the product of the two node factors is \(144u^4\), a square, so the pattern weight remains \(\langle1\rangle\). Proposition 2.8 gives the same aggregate trace \(\langle3\rangle+h\) for the corrected cubic algebra. This is positive evidence that the final formula may survive an appropriate repair.\par
Retain the coefficient-dependent Kummer algebra, or justify the normalization together with its descent over \(F\); then redo the pattern-to-curve identification and the successive transfers. Neither the special normalized case of (33) nor equality of the final trace in this example proves the general descent step.\par
Problem 3 (Minor: unrestricted tropicalization need not preserve genus).  The paragraph immediately before Remark 3.5 and the remark itself, p. 15, claim genus preservation without the later genericity restrictions. Over \(\mathbb C\{\{t\}\}\), the constant smooth cubic \[X^3+Y^3+Z^3=0\] has geometric genus one. Its embedded plane tropicalization is the corner locus of \(\max\{0,3x,3y\}\), a three-ray tree with first Betti number zero. Thus this assertion needs a degeneration hypothesis or a tropical genus that includes vertex genera. The example does not satisfy the simple/maximal-rank situation of Definition 3.10 used for the correspondence theorem. State the appropriate restricted genus comparison rather than applying the unconditional sentence. No contradiction to the generic characteristic-zero correspondence is obtained from this example.\par
Problem 4 (Minor: the toric variety requires the normal fan).  Section 2.4, p. 12, constructs \(\operatorname{Tor}_k(\Delta)\) from cones over the proper faces of \(\Delta\), then identifies the standard degree-\(d\) triangle with \(\mathbb P^2\). Already for \(\Delta=\operatorname{Conv}\{0,e_1,e_2\}\), the stated cones form the first-quadrant cone and its faces; their toric variety is \(\mathbb A^2\), not \(\mathbb P^2\). Use the normal fan in the dual lattice, or an explicitly appropriate polar construction. The later Newton-polygon and projective-boundary arguments plainly require the usual projective toric surface; the evidence here isolates a false background construction, not an additional failure of the transfer identities.\par
\section*{Checks supporting the formula and explicit limits}
For a separable algebra \(F[x]/(x^m-D)\), the trace matrix in the power basis has entries supported when the exponent is divisible by \(m\). Pairing \(x^i\) with \(x^{m-i}\) yields the hyperbolic blocks in Proposition 2.8; the remaining entries give its stated odd/even rank-one terms. Eight exact power-basis examples were computed. The characteristic-prime condition is essential.\par
The square-class invariance of a nodal Hessian under a coordinate change was checked: the determinant is multiplied by the square of the Jacobian determinant; multiplying the defining equation by a unit likewise contributes a square at a node. The parallelogram Hessian calculation gives \((A-B)^2\) when the intersections are transverse. The root-product calculation for triangles was followed with its nonzero-denominator restrictions.\par
In Section 6.2 the odd-edge calculation is explicitly performed in the Witt ring, so omitted hyperbolic summands in that calculation are not errors. Every interior extended edge enters twice, while boundary edges have length one. Together with the rank product, this produces the stated anisotropic sign. An even edge yields a hyperbolic transfer, again conditional on valid étale coefficient algebras. The del Pezzo corollary additionally uses the cited invariance and specialization results; rank and real signature refer to the universal integer combination of \(\langle1\rangle,\langle-1\rangle\), rather than a signature operation on an arbitrary field.\par
Actual inspection of p. 5 shows a closure bar in (4), lost by text extraction. The apparent missing closure was therefore rejected as a finding. The cardinality-only wording in the proof of Theorem 4.1 does not by itself prove bijectivity, but the present review does not promote that observation into another confirmed defect without a separate counterexample or a full audit of the cited reconstruction theorem.\par
The review did not reprove all log/toric deformation results, global regularity criteria, finite-field existence statements, or the external del Pezzo moduli-space invariance and specialization theorems. It does not certify a repaired positive-characteristic bound or all coefficient twists. Searches covered the exact arXiv version, publisher and both authors’ research pages; no separate supplement or author code was located. No author code was executed.\par
\section*{Scope and limitations}
The original accented names, frozen identifier and year are retained. The source, exact computations, TeX, existing compiler, dependency payload, PDF and inspected images are hash-bound in the accompanying records. This revision incorporates focused independent acceptance of the existing findings, including the explicit compatibility check above. The original review’s coverage and stated limitations are retained. Scientific adjudication remains pending; human reviews are empty.\par
\end{document}
